%-------------------------------------------------------------------------------
\section{Experimental Evaluation}
\label{sec:eval}
%-------------------------------------------------------------------------------

This section runs experiments with the methods and framework above and draws
quantitative conclusions. Besides the internal layered metrics
(the end-to-end result snapshot below and the lint-coverage analysis of
\S\ref{subsec:coverage}), it provides two
mutually independent external validations, each checking one half of the
framework: \S\ref{subsec:lintability-ext} externally validates rule
extraction and lintability determination against the zlint maintainers' manual
gold standard, and \S\ref{subsec:certdetect} externally validates code
generation and synonymy judgment by execution on real certificates.

\subsection*{Experimental Setup and Result Snapshot}

Evaluation is performed on full-corpus re-extraction results of the two standard
sources RFC~5280 and CABF BR. The LLM used in the experiments is GPT-5.4.

\begin{table*}[tb]
\centering
\footnotesize
\caption{Current end-to-end result snapshot (RFC~5280 $+$ CABF BR).}
\label{tab:snapshot}
\begin{tabular}{@{}p{4.4cm}r p{10.4cm}@{}}
\toprule
Metric & Count & Scope / notes \\
\midrule
Total recalled rules & 2080 & RFC~5280 640 $+$ CABF BR 1440 \\
lintable (single-cert observable) & 248 &
  current re-extraction / re-adjudication: CABF 170 $+$ RFC~5280 78; the negative
  gate excludes process, cross-certificate, runtime, external-semantic, and
  MAY/OPTIONAL classes; CRL document rules are outside this scope \\
zlint coverage (full; existing cognate lint) & 158 &
  the rule is fully implemented by some cognate zlint certificate lint; residual
  cross-standard inheritance gaps are analyzed in \S\ref{subsec:crossstd} \\
uncovered lintable (codegen domain) & 90 &
  enters the current codegen domain ($248 - 158$); CABF 64 $+$ RFC~5280 26, the
  target set of downstream code generation $\phi_G$ \\
can generate compilable lint & 90 &
  deterministic path 87 $+$ LLM fallback path 3, both gated by ``render and
  pass go build''; generation rate $= 90/90 = 100.0\%$ \\
cannot generate (no tree / not compiling / abstain) & 0 &
  no generation failure in the current codegen domain \\
final-shipping strict synonymous (\CodeEqSpec) & 90 &
  the final in-tree zlint lint (\texttt{CheckApplies} $+$ \texttt{Execute} $+$
  severity/date metadata) judged against the original rule text and context;
  strict synonymy rate $= 90/90 = 100.0\%$ \\
final-shipping not synonymous (does not express) & 0 &
  no final-shipping residual in the current codegen domain \\
emitted lints' atom templates touched & 52 &
  28 GENERIC $+$ 24 NON\_GENERIC of the 136 registered atoms; by rule, 53
  GENERIC-only / 5 mixed / 32 NON\_GENERIC-only \\
\emph{(collateral) certificate-level execution verification \CodeEqIR} &
  \emph{37} &
  \emph{model-independent IR-level faithfulness, not a synonymy numerator; the
  paper-facing synonymy rate uses only the final emitted code versus the source
  rule} \\
\bottomrule
\end{tabular}
\end{table*}

Under the current codegen domain there is no non-generable compilable lint and no
final-shipping non-synonymous residual (both counts are 0). The 37 rules
certified \CodeEqIR{} by certificate-level execution verification are reported
only as collateral IR-level evidence: \CodeEqIR{} is model-independent and
deterministic, yet it does not by itself imply \CodeEqSpec, so the paper-facing
synonymy rate is taken solely from the final emitted code judged against the
source rule. Any reduction of the denominator therefore comes from re-extraction
and lintability re-adjudication rather than from local metric filters; e.g.,
R31068 (RFC~5280~\S4.2.1.6) was re-adjudicated as not lintable because a single
certificate exposes only the chosen GeneralName tag and its value, not the
issuer's pre-encoding intent.

The first-layer quality signal is G1 conservation (Theorem~\ref{thm:recall},
\S\ref{subsec:targets}): the total number of keyword-recalled rules should
neither increase nor decrease in any downstream classification step. Under the
current snapshot, conservation holds strictly:
\[
  \underbrace{2080}_{\text{recall}} \;=\; \underbrace{529}_{\text{noise}} +
  \underbrace{1551}_{\text{genuine}}, \quad
  \underbrace{1551}_{\text{genuine}} \;=\; \underbrace{248}_{\text{lintable}} +
  \underbrace{1303}_{\text{not lintable}}.
\]
Both identities hold term by term, i.e., the G1 residual $\Lrec = 0$ and the top
two layers of conservation are closed (deterministically recomputable). The 248
lintable rules split by standard source
into 170 CABF and 78 RFC~5280; of these zlint fully covers 158, leaving
90---the domain of downstream code generation $\phi_G$.

\subsection{Lint Coverage Analysis}
\label{subsec:coverage}

Answering ``how many lintable rules do existing tools cover'' hinges on using
the right criterion. We use ``whether the lintable rule is genuinely implemented
by some cognate zlint lint'' as the criterion: for each rule we retrieve
candidate zlint lints, then compare field by field (subject / obligation /
predicate / constraint), giving a three-level verdict full (fully implemented) /
partial (partially implemented) / none (not implemented).

Fig.~\ref{alg:coverage} formalizes this determination: summarize a lint
into a reverse IR, then align field by field; candidate retrieval is by
source/section (Stage-1: RFC section numbers are stable and narrowed by prefix,
while CABF section numbers drift across versions so all CABF lints are taken),
with an added ``wrong-field'' consistency gate (Stage-3, only downgrades), and
coverage counts only full (Stage-4; partial/none go to $\phi_G$).

% Algorithm 1 -- zlint coverage-alignment determination for lintable normative rules.
% Included from sections/08_evaluation.tex via % Algorithm 1 -- zlint coverage-alignment determination for lintable normative rules.
% Included from sections/08_evaluation.tex via % Algorithm 1 -- zlint coverage-alignment determination for lintable normative rules.
% Included from sections/08_evaluation.tex via \input{algorithms/alg_coverage}.
%
% Typeset single-column (not figure*) at \scriptsize: the algorithmic lines are
% short, so at full width the right half of most lines was blank and the float was
% billed two columns of page area. Long \Comment texts are placed on their own
% \Statex lines because \Comment is right-aligned and cannot wrap. No step, symbol,
% or comment wording was changed.
\begin{figure}[tb]
\makeatletter
\@ifundefined{c@algorithm}{%
  \newcounter{algorithm}%
  \renewcommand{\thealgorithm}{\arabic{algorithm}}%
}{}%
\makeatother
\refstepcounter{algorithm}
\label{alg:coverage}
\noindent\rule{\linewidth}{0.6pt}
\par\smallskip
\noindent\textbf{Algorithm~\thealgorithm.} Algorithmic listing for zlint
coverage-alignment determination for individual normative rules.
\par\smallskip
\noindent\rule{\linewidth}{0.4pt}
\par\smallskip
\scriptsize
\begin{algorithmic}[1]
\Statex \textbf{Input:} lintable normative-rule set $B=\{b_1,\dots,b_m\}$, where each normative rule has
  IR four-tuple $\langle$subject, obligation, predicate, constraint$\rangle$,
  source, and section; zlint lint set $L=\{\ell_1,\dots,\ell_n\}$, where each lint has
  description/citation/severity metadata and Pass/Error tests; offline code
  summarizer $M_{\mathrm{summ}}$, and field-level judge $M_{\mathrm{judge}}$
  with temperature $=0$
\Statex \textbf{Output:} per-normative-rule coverage verdict $\mathrm{verdict}(b)\in\{\text{full, partial,
  none}\}$, covered set $\RLcov$, and uncovered code-generation domain $\RLunc$
\Statex // Offline cached Stage-0: summarize each lint into a reverse
  IR
\ForAll{$\ell \in L$}
  \State $\ell.\mathit{ir} \gets M_{\mathrm{summ}}(\ell.\mathrm{src},
    \ell.\mathrm{tests}, \ell.\mathrm{meta})$
  \Statex \Comment{four-tuple $+$ code\_summary}
\EndFor
\State pre-split candidate pools by Source: $L_{\mathrm{RFC}}\gets$ RFC lints,\;
  $L_{\mathrm{CABF}}\gets$ CABF lints
\ForAll{$b \in B$}
  \State // Stage-1: deterministic candidate retrieval
  \If{$b.\mathrm{source}=\mathrm{RFC}$}
    \State $C(b)\gets\{\ell\in L_{\mathrm{RFC}}:
      \textsc{SectionPrefixMatch}(b.\mathrm{section},\ell.\mathrm{citedSections})\}$
  \ElsIf{$b.\mathrm{source}=\mathrm{CABF}$}
    \State $C(b)\gets L_{\mathrm{CABF}}$ \Comment{CABF sections drift $\Rightarrow$ take all}
  \Else
    \State $C(b)\gets\emptyset$
  \EndIf
  \State // Stage-2: field-level judge, compare four-tuple, best across candidates
  \State $(v,\ell^{*})\gets M_{\mathrm{judge}}(b.\mathit{ir}, C(b))$
  \Statex \Comment{$v\in\{\text{full},\text{partial},\text{none}\}$}
  \State // Stage-3: reject positive matches with conflicting concrete subject families
  \If{$v\in\{\text{full},\text{partial}\} \wedge
      \mathrm{Family}(b.\mathrm{subject})\neq\emptyset \wedge
      \mathrm{Family}(\ell^{*}.\mathrm{subject})\neq\emptyset \wedge
      \mathrm{Family}(b.\mathrm{subject})\neq\mathrm{Family}(\ell^{*}.\mathrm{subject})$}
    \State $v\gets\text{none}$
    \Statex \Comment{different concrete families $\Rightarrow$ wrong-field false positive}
  \EndIf
  \State $\mathrm{verdict}(b)\gets v$
\EndFor
\State // Stage-4: deterministic aggregation, coverage counts full only
\State $\RLcov\gets\{b\in B:\mathrm{verdict}(b)=\text{full}\}$;\quad
  $\RLunc\gets B\setminus \RLcov$
\State \textbf{return} $(\{(b,\mathrm{verdict}(b)):b\in B\},\ \RLcov,\ \RLunc)$
\end{algorithmic}
\par\smallskip
\noindent\rule{\linewidth}{0.6pt}
\end{figure}
.
%
% Typeset single-column (not figure*) at \scriptsize: the algorithmic lines are
% short, so at full width the right half of most lines was blank and the float was
% billed two columns of page area. Long \Comment texts are placed on their own
% \Statex lines because \Comment is right-aligned and cannot wrap. No step, symbol,
% or comment wording was changed.
\begin{figure}[tb]
\makeatletter
\@ifundefined{c@algorithm}{%
  \newcounter{algorithm}%
  \renewcommand{\thealgorithm}{\arabic{algorithm}}%
}{}%
\makeatother
\refstepcounter{algorithm}
\label{alg:coverage}
\noindent\rule{\linewidth}{0.6pt}
\par\smallskip
\noindent\textbf{Algorithm~\thealgorithm.} Algorithmic listing for zlint
coverage-alignment determination for individual normative rules.
\par\smallskip
\noindent\rule{\linewidth}{0.4pt}
\par\smallskip
\scriptsize
\begin{algorithmic}[1]
\Statex \textbf{Input:} lintable normative-rule set $B=\{b_1,\dots,b_m\}$, where each normative rule has
  IR four-tuple $\langle$subject, obligation, predicate, constraint$\rangle$,
  source, and section; zlint lint set $L=\{\ell_1,\dots,\ell_n\}$, where each lint has
  description/citation/severity metadata and Pass/Error tests; offline code
  summarizer $M_{\mathrm{summ}}$, and field-level judge $M_{\mathrm{judge}}$
  with temperature $=0$
\Statex \textbf{Output:} per-normative-rule coverage verdict $\mathrm{verdict}(b)\in\{\text{full, partial,
  none}\}$, covered set $\RLcov$, and uncovered code-generation domain $\RLunc$
\Statex // Offline cached Stage-0: summarize each lint into a reverse
  IR
\ForAll{$\ell \in L$}
  \State $\ell.\mathit{ir} \gets M_{\mathrm{summ}}(\ell.\mathrm{src},
    \ell.\mathrm{tests}, \ell.\mathrm{meta})$
  \Statex \Comment{four-tuple $+$ code\_summary}
\EndFor
\State pre-split candidate pools by Source: $L_{\mathrm{RFC}}\gets$ RFC lints,\;
  $L_{\mathrm{CABF}}\gets$ CABF lints
\ForAll{$b \in B$}
  \State // Stage-1: deterministic candidate retrieval
  \If{$b.\mathrm{source}=\mathrm{RFC}$}
    \State $C(b)\gets\{\ell\in L_{\mathrm{RFC}}:
      \textsc{SectionPrefixMatch}(b.\mathrm{section},\ell.\mathrm{citedSections})\}$
  \ElsIf{$b.\mathrm{source}=\mathrm{CABF}$}
    \State $C(b)\gets L_{\mathrm{CABF}}$ \Comment{CABF sections drift $\Rightarrow$ take all}
  \Else
    \State $C(b)\gets\emptyset$
  \EndIf
  \State // Stage-2: field-level judge, compare four-tuple, best across candidates
  \State $(v,\ell^{*})\gets M_{\mathrm{judge}}(b.\mathit{ir}, C(b))$
  \Statex \Comment{$v\in\{\text{full},\text{partial},\text{none}\}$}
  \State // Stage-3: reject positive matches with conflicting concrete subject families
  \If{$v\in\{\text{full},\text{partial}\} \wedge
      \mathrm{Family}(b.\mathrm{subject})\neq\emptyset \wedge
      \mathrm{Family}(\ell^{*}.\mathrm{subject})\neq\emptyset \wedge
      \mathrm{Family}(b.\mathrm{subject})\neq\mathrm{Family}(\ell^{*}.\mathrm{subject})$}
    \State $v\gets\text{none}$
    \Statex \Comment{different concrete families $\Rightarrow$ wrong-field false positive}
  \EndIf
  \State $\mathrm{verdict}(b)\gets v$
\EndFor
\State // Stage-4: deterministic aggregation, coverage counts full only
\State $\RLcov\gets\{b\in B:\mathrm{verdict}(b)=\text{full}\}$;\quad
  $\RLunc\gets B\setminus \RLcov$
\State \textbf{return} $(\{(b,\mathrm{verdict}(b)):b\in B\},\ \RLcov,\ \RLunc)$
\end{algorithmic}
\par\smallskip
\noindent\rule{\linewidth}{0.6pt}
\end{figure}
.
%
% Typeset single-column (not figure*) at \scriptsize: the algorithmic lines are
% short, so at full width the right half of most lines was blank and the float was
% billed two columns of page area. Long \Comment texts are placed on their own
% \Statex lines because \Comment is right-aligned and cannot wrap. No step, symbol,
% or comment wording was changed.
\begin{figure}[tb]
\makeatletter
\@ifundefined{c@algorithm}{%
  \newcounter{algorithm}%
  \renewcommand{\thealgorithm}{\arabic{algorithm}}%
}{}%
\makeatother
\refstepcounter{algorithm}
\label{alg:coverage}
\noindent\rule{\linewidth}{0.6pt}
\par\smallskip
\noindent\textbf{Algorithm~\thealgorithm.} Algorithmic listing for zlint
coverage-alignment determination for individual normative rules.
\par\smallskip
\noindent\rule{\linewidth}{0.4pt}
\par\smallskip
\scriptsize
\begin{algorithmic}[1]
\Statex \textbf{Input:} lintable normative-rule set $B=\{b_1,\dots,b_m\}$, where each normative rule has
  IR four-tuple $\langle$subject, obligation, predicate, constraint$\rangle$,
  source, and section; zlint lint set $L=\{\ell_1,\dots,\ell_n\}$, where each lint has
  description/citation/severity metadata and Pass/Error tests; offline code
  summarizer $M_{\mathrm{summ}}$, and field-level judge $M_{\mathrm{judge}}$
  with temperature $=0$
\Statex \textbf{Output:} per-normative-rule coverage verdict $\mathrm{verdict}(b)\in\{\text{full, partial,
  none}\}$, covered set $\RLcov$, and uncovered code-generation domain $\RLunc$
\Statex // Offline cached Stage-0: summarize each lint into a reverse
  IR
\ForAll{$\ell \in L$}
  \State $\ell.\mathit{ir} \gets M_{\mathrm{summ}}(\ell.\mathrm{src},
    \ell.\mathrm{tests}, \ell.\mathrm{meta})$
  \Statex \Comment{four-tuple $+$ code\_summary}
\EndFor
\State pre-split candidate pools by Source: $L_{\mathrm{RFC}}\gets$ RFC lints,\;
  $L_{\mathrm{CABF}}\gets$ CABF lints
\ForAll{$b \in B$}
  \State // Stage-1: deterministic candidate retrieval
  \If{$b.\mathrm{source}=\mathrm{RFC}$}
    \State $C(b)\gets\{\ell\in L_{\mathrm{RFC}}:
      \textsc{SectionPrefixMatch}(b.\mathrm{section},\ell.\mathrm{citedSections})\}$
  \ElsIf{$b.\mathrm{source}=\mathrm{CABF}$}
    \State $C(b)\gets L_{\mathrm{CABF}}$ \Comment{CABF sections drift $\Rightarrow$ take all}
  \Else
    \State $C(b)\gets\emptyset$
  \EndIf
  \State // Stage-2: field-level judge, compare four-tuple, best across candidates
  \State $(v,\ell^{*})\gets M_{\mathrm{judge}}(b.\mathit{ir}, C(b))$
  \Statex \Comment{$v\in\{\text{full},\text{partial},\text{none}\}$}
  \State // Stage-3: reject positive matches with conflicting concrete subject families
  \If{$v\in\{\text{full},\text{partial}\} \wedge
      \mathrm{Family}(b.\mathrm{subject})\neq\emptyset \wedge
      \mathrm{Family}(\ell^{*}.\mathrm{subject})\neq\emptyset \wedge
      \mathrm{Family}(b.\mathrm{subject})\neq\mathrm{Family}(\ell^{*}.\mathrm{subject})$}
    \State $v\gets\text{none}$
    \Statex \Comment{different concrete families $\Rightarrow$ wrong-field false positive}
  \EndIf
  \State $\mathrm{verdict}(b)\gets v$
\EndFor
\State // Stage-4: deterministic aggregation, coverage counts full only
\State $\RLcov\gets\{b\in B:\mathrm{verdict}(b)=\text{full}\}$;\quad
  $\RLunc\gets B\setminus \RLcov$
\State \textbf{return} $(\{(b,\mathrm{verdict}(b)):b\in B\},\ \RLcov,\ \RLunc)$
\end{algorithmic}
\par\smallskip
\noindent\rule{\linewidth}{0.6pt}
\end{figure}


Here \textsc{SectionPrefixMatch} is true when the rule's section number shares a
common prefix with a lint's cited section number; $\mathrm{Family}(\cdot)$ maps a
subject path to a concrete field family, returning $\emptyset$ for a
vague/unresolved subject rather than downgrading (so genuine matches on coarse
subjects are preserved). The algorithm calls the LLM only in Stage-0 and Stage-2
(the judge at temperature $=0$, its prompt containing three kinds of
positive/negative examples: direction reversal, field misplacement, and
constraint-type confusion); summaries are cached offline, the full CABF
candidate set is sent to the judge in batches, and the rest is deterministic.

Of the 248 lintable rules, zlint fully covers 158 and leaves 90
uncovered under this coverage procedure (Table~\ref{tab:coverage}); the 90
uncovered are the domain of code generation $\phi_G$.

\begin{table}[tb]
\centering
\small
\caption{zlint cognate lint counts and their coverage of the 248 lintable
rules, by standard source.}
\label{tab:coverage}
\begin{tabular}{@{}lrrr@{}}
\toprule
Item & CABF & RFC~5280 & Total \\
\midrule
\emph{zlint cognate lints, total (ref.)} & \emph{172} & \emph{123} & \emph{295} \\
\quad\emph{— of which cert.\ lints} & \emph{166} & \emph{115} & \emph{281} \\
\quad\emph{— of which CRL lints} & \emph{6} & \emph{8} & \emph{14} \\
full (full coverage) & 106 & 52 & 158 \\
uncovered (codegen domain) & 64 & 26 & 90 \\
lintable total & 170 & 78 & 248 \\
\bottomrule
\end{tabular}
\end{table}

The first three rows (italic) are lint counts computed directly from the
project's built-in v3 source by zlint's Source metadata field (unit:
lints; the 14 CRL lints are identified via RegisterRevocationListLint
and fall outside our single-cert scope); the remaining rows are the coverage
verdicts of our lintable rules (unit: rules). ``full'' counts ``whether one of
our rules is fully implemented by some cognate zlint lint''; it does not form a
simple ratio with the total zlint lint count, mainly because one lint can hit
multiple rules.

\subsection{Lintability External Validation: Rule Extraction and Lintability
Determination}
\label{subsec:lintability-ext}

The first half of the framework---rule extraction and lintability
determination---needs an independent external reference. For this we bring in
the CABF BR mapping table published by zlint
maintainers~\cite{zlint-mapping-sheet}: this table is manually annotated rule by rule
by zlint experts as to ``whether a BR requirement can be expressed by a static
lint,'' unaware of our extractor, IR, and lintability decider. Each Yes/No label
in the table corresponds exactly to our joint output ``extract this rule $+$
decide whether it is lintable'' for the same rule; agreement indicates that our
rule extraction and lintability determination match human experts. The table
covers two versions, BR~1.4.8~\cite{cabf-br148-md} and BR~2.0.2~\cite{cabf-br202-md}, compared version-for-version against
our CABF-Server-1.4.8 backend (422 rules) and CABF-Server-2.0.2
backend (1024 rules) extracted by the same two-layer
pipeline.

\begin{table}[tb]
\centering
\scriptsize
\setlength{\tabcolsep}{3pt}
\caption{Version-for-version external validation on BR~1.4.8 / BR~2.0.2.}
\label{tab:extval}
\begin{tabular}{@{}lrrrrr@{}}
\toprule
Version & Sheet rows (Yes/No) & Match cov. & Agreement &
  Cohen's $\kappa$ & F1 \\
\midrule
BR~1.4.8 & 122 (75/47) & 86.9\% & 91.5\% & 0.823 & 0.930 \\
BR~2.0.2 & 250 (17/233) & 86.8\% & 96.8\% & 0.616 & 0.632 \\
\bottomrule
\end{tabular}
\end{table}

A small number of sheet rows fall outside our scope---title-case or lowercase descriptors such as ``Required if'', ``Deprecated'', or lowercase ``shall'' that RFC~2119 processing correctly skips, and, on the BR~2.0.2 sheet, cross-references or chapter anchors with no matching backend clause---so match coverage is near 87\% on both versions. On BR~2.0.2 the matched positive class is very small (12 of 217 matched pairs, 5.5\%), so a naive all-No baseline already attains 94.5\% agreement; thus $\kappa = 0.616$ is the more meaningful figure, and the positive-class metrics ($P = 0.857$, $R = 0.50$, $n = 12$) carry correspondingly wide confidence intervals. BR~1.4.8 has more balanced classes (69 of 106 matched pairs positive, 65\%) and yields $\kappa = 0.823$ with positive-class $F_1 = 0.930$ and zero false positives, a tighter estimate of system quality.

This external validation checks the first half of the system. On
BR~1.4.8, $\kappa = 0.823$ and Yes recall 87.0\% indicate that our rule
extraction (identifying which of the 422/1024 backend items are genuine
normative requirements) and lintability determination (the four-condition
conjunction, \S\ref{sec:lintability}) agree strongly with zlint experts'
independent manual judgment; the 50.0\% Yes recall on BR~2.0.2 indicates a clear
false-negative margin remaining in the newer BR version. This validation does
not check downstream code generation and synonymy; those are checked
additionally on real certificates by the certificate detection of
\S\ref{subsec:certdetect}.

\subsection{Certificate Detection}
\label{subsec:certdetect}

Injecting the 90 final-shipping-strict lints into the zlint binary and scanning
certificates checks code generation and synonymy judgment. We count a hit as a
valid finding only when three conditions hold: the final zlint code's semantics
have been checked back to the source text, the certificate yields a non-pass
result, and an independent structural audit re-derives the same defect from the
DER / \texttt{openssl}; on external corpora we additionally require that no native
zlint lint flags the same certificate, so that existing coverage is not
miscredited to a generated lint. The unit is a (certificate $\times$ lint)
finding. The corpus first uses zlint's own testdata as a gate rather than an
ecosystem estimate: on 1128 parsable certificates the 90 lints raise 2627 raw
cicasgen\_ hits, of which 299 are independently confirmed (251 Error $+$
48 Warn) and 2328 are \texttt{NOCHECK} and hence not claimed.

\begin{table*}[tb]
\centering
\scriptsize
\setlength{\tabcolsep}{4pt}
\caption{Strictly confirmed findings across the three certificate corpora
(highest severity first; 90 shipped lints). ``ext.\ merged'' is the Tranco/CT
union after cross-corpus de-duplication.}
\label{tab:certfind}
\begin{tabular}{@{}l p{5.9cm} r r r r p{3.3cm}@{}}
\toprule
Sev. & Finding type / governing obligation & testdata & Tranco & CT &
  ext.\ merged & Evidence boundary \\
\midrule
Error & Subscriber-cert Certificate Policies must contain exactly one CABF
  Reserved Certificate Policy Identifier (BR~\S7.1.2.7.9) & 113 & 1 & 0 & 1 &
  one Tranco leaf lacks the reserved policy OID; no native zlint overlap \\
Error & AKI must not contain authorityCertSerialNumber / authorityCertIssuer
  (BR~\S7.1.2.11.1) & 32 & 2 & 0 & 2 & one Tranco CA cert violates both AKI
  subfields; no native zlint overlap \\
Error & Subscriber-cert Certificate Policies must not contain anyPolicy
  (BR~\S7.1.2.7.9) & 26 & 0 & 0 & 0 & confirmed only in testdata fixtures \\
Error & IV/OV subject-profile constraints on surname, givenName, and
  localityName / stateOrProvinceName (BR~\S7.1.2.7.3--\S7.1.2.7.4) & 57 & 0 & 0 &
  0 & testdata fixtures only \\
Error & Other structural constraints: SAN critical, signatureAlgorithm matches
  tbsCertificate.signature, subscriber pathLenConstraint forbidden, empty
  Certificate Policies, v1 uniqueID forbidden (RFC~5280 / BR) & 23 & 0 & 0 & 0 &
  testdata fixtures only \\
Warn & A certificate carrying only basic fields SHOULD be v1
  (RFC~5280~\S4.1.2.1) & 47 & 0 & 0 & 0 & advisory, reported as lint.Warn \\
Warn & Root CA Certificate SHOULD NOT contain CRL Distribution Points
  (BR~\S7.1.2.11.2) & 1 & 3 & 1 & 3 & advisory profile warning; the CT finding
  duplicates the same Root CA as one Tranco finding \\
\bottomrule
\end{tabular}
\end{table*}

The testdata column sums to 299 confirmed fixture findings (251 Error $+$ 48
Warn); Tranco contributes 6 and CT 1, i.e., 7 findings by source. We read
severity conservatively: the Certificate-Policies and AKI violations are
Error-level profile faults whose direct cryptographic consequence is indirect,
while the Root CA CRLDP finding is a lower, Warn-level advisory.

\textbf{Supplementary scan on real certificates.} To test whether these findings
exist only in hand-constructed testdata, we scan two external corpora: the
TLS-handshake certificate chains of the Tranco Top~1M domains (47,791
deduplicated PEMs) and a recent sample of CT logs (63,327 deduplicated PEMs). The
raw cicasgen\_ hit counts are large---44,020 on Tranco and 57,558 on CT---but
cannot be taken as defects: if a generated lint loses profile scope, a
conditional clause, or an effective-date constraint, an independent audit may
still only prove that the certificate satisfies the trigger of the wrong lint. We
therefore treat the external corpora only as negative validation. Under the
strict no-upstream audit these raw hits reduce to 6 confirmed on Tranco and 1 on
CT; summed by source that is 7 findings, and after cross-corpus de-duplication it
is 6 (certificate $\times$ lint) findings over 5 unique certificates, because one
Root CA CRLDP warning is confirmed on the same Comodo ``AAA Certificate
Services'' root in both corpora.

Raw hit rates therefore serve only to surface audit candidates, never as
defect-count conclusions. This supports the negative message of
Table~\ref{tab:snapshot}: even when a lint compiles, executes, and fires on a
certificate, one must still re-check the source-text scope, its conditional
clauses, and rule effective dates before calling it a defect. Concretely,
several high-frequency cicasgen\_ hits are excluded on exactly these
grounds---an OV/IV localityName requirement that dropped its
``stateOrProvinceName absent'' precondition (all such certs carry
stateOrProvinceName); an organizationalUnitName or AKI-subfield prohibition
applied to certificates issued before the requirement's effective date; the
RFC~5280 A.1 ``version MUST be v2'' fragment mis-extracted as a profile rule that
flags normal v3 certificates; and SubCA anyPolicy / NameConstraints rules that
approximate profile scope with IsSubCA alone. Hence the raw hit rate must not be
read as the real compliance-defect rate.
